%%%%%%%%%%  scan idx 114  |  folio 101  %%%%%%%%%%
\SGAPrintPage{101}%
\uline{is fully faithful and has essential image} \(\Dpx(\underline{C})_{\underline{C}_0\text{-}\mathrm{coh}}\).

\uline{Proof.} a) Let \(F\in\Ob\Dpx(\underline{C})_{\underline{C}_0\text{-}\mathrm{coh}}\), and suppose that an \(i\)-quasi-isomorphism (of \(\Cpx(\underline{C})\)) \(u_i:E_i\longrightarrow F\) has been constructed with \(E_i\) strictly pseudo-coherent (this is possible for \(i\) sufficiently large with \(E_i=0\)). Let \(G_i\) be the cone of \(u_i\). It follows from (1.4.2) and (2.1.1) that \(\Ho^{i-1}(G_i)\) is of \(\underline{C}_0\)-finite type. Thus, by (1.4.4), there exists a commutative triangle of \(\Cpx(\underline{C})\)
\[
\begin{tikzcd}
E_i \arrow[r,"v_{i-1}"] \arrow[d,"u_i"'] & E_{i-1} \arrow[dl,"u_{i-1}"]\\
F & ,
\end{tikzcd}
\]
where \(E_{i-1}\) is strictly pseudo-coherent, \(u_{i-1}\) an \((i-1)\)-quasi-isomorphism, and \(v_{i-1}\) a monomorphism inducing the identity in degree \(\geq i\). By passage to the limit (cf.\ proof of (1.4.5)), we deduce a quasi-isomorphism \(u:E\to F\) with \(E\) strictly pseudo-coherent.

b) Apply the criterion ([V] chap I \S\ 2 4.2 b)(i)) and a).

Observe that the proof of a) gives the following more precise result, which generalizes (1.4.5):

\sgaStableTarget{sga6.e01.prop.2-7-1}\textbf{Proposition 2.7.1.} \uline{Under the hypotheses of (2.7), let} \(u:E\to F\) \uline{be an} \(n\)\uline{-quasi-isomorphism of} \(\Cpx(\underline{C})\)\uline{, with} \(F\) \uline{pseudo-coherent and} \(E\) \uline{strictly pseudo-coherent. There then exists a commutative triangle of} \(\Cpx(\underline{C})\)
\[
\begin{tikzcd}
E \arrow[r,"v"] \arrow[d,"u"'] & E' \arrow[dl,"u'"]\\
F & ,
\end{tikzcd}
\]
\uline{where} \(E'\) \uline{is strictly pseudo-coherent,} \(u'\) \uline{a quasi-isomorphism, and} \(v\) \uline{a}
%%%%%%%%%%  scan idx 115  |  folio 102  %%%%%%%%%%
\SGAPrintPage{102}%
\uline{monomorphism inducing the identity in degree} \(\geq n\).

\sgaStableTarget{sga6.e01.def.2-8}\textbf{Definition 2.8.} \uline{Let} \(n\in\mathbb N\) \uline{and} \(F\in\Ob\underline{C}_S\)\uline{. We shall say that} \(F\) \uline{is of} \(\underline{n}\)-\(\underline{C}_0\)-finite presentation \uline{when the set of} \(X\) \uline{over} \(S\) \uline{such that there exists an exact sequence}
\[
  E^{-n}\longrightarrow\ldots\longrightarrow E^0\longrightarrow F_X\longrightarrow 0,
\]
\uline{with} \(E^i\in\Ob\underline{C}_{0X}\) \uline{for} \(-n\leq i\leq 0\)\uline{, is a refinement of} \(S\)\uline{. We shall say that} \(F\) \uline{is of} \(\infty\)-\(\underline{C}_0\)-finite presentation \uline{when} \(F\) \uline{is of} \(n\)\uline{-}\(\underline{C}_0\)\uline{-finite presentation for every} \(n\in\mathbb N\).

Of course, we shall sometimes simply say \(n\)-finite presentation when no ambiguity concerning \(\underline{C}_0\) is possible.

We shall say ``of finite presentation'' for ``of 1-finite presentation'', and ``of finite type'' for ``of 0-finite presentation'', the latter notion coinciding with the one already introduced in (1.2).

\sgaStableTarget{sga6.e01.prop.2-9}\textbf{Proposition 2.9.} \uline{Let} \(n\in\mathbb N\) \uline{and} \(F\in\Ob\underline{C}_S\)\uline{. Equivalent conditions}:

(i) \(F\) \uline{is of} \(n\)\uline{-finite presentation (resp.\ of} \(\infty\)\uline{-finite presentation)};

(ii) \(F\) \uline{is} \((-n)\)\uline{-pseudo-coherent (resp.\ pseudo-coherent)}.

\uline{In the case where} \(\underline{S}\) \uline{is a punctual site,} \(F\) \uline{is pseudo-coherent if and only if there exists an exact sequence}
\[
  \ldots\longrightarrow E^i\longrightarrow\ldots\longrightarrow E^0\longrightarrow F\longrightarrow 0,
\]
\uline{with} \(E^i\in\Ob\underline{C}_0\) \uline{for every} \(i\).

\uline{Proof.} For the equivalence \((\mathrm{i})\Leftrightarrow(\mathrm{ii})\), it clearly suffices to treat the case without the ``resp.'' alternatives. The implication \((\mathrm{i})\Rightarrow(\mathrm{ii})\) follows immediately from the definitions, and
%%%%%%%%%%  scan idx 116  |  folio 103  %%%%%%%%%%
\SGAPrintPage{103}%
the implication \((\mathrm{ii})\Rightarrow(\mathrm{i})\) follows at once from (2.4): the arrow \(0\to F\) is a 1-quasi-isomorphism; we deduce (after localizing) a quasi-isomorphism \(E\longrightarrow F\) with \(E\) strictly \((-n)\)-pseudo-coherent and \(E^i=0\) for \(i\geq 1\), i.e.\ an \(n\)-finite presentation of \(F\). The second assertion is proved in a similar way (use (2.7.1)).

\sgaStableTarget{sga6.e01.prop.2-10}\textbf{Proposition 2.10.} a) \uline{Let} \(m,n\in\mathbb Z\)\uline{, with} \(n\leq m\)\uline{. When} \(F\in\Ob\Dpx^-(\underline{C}_S)\) \uline{is} \(n\)\uline{-pseudo-coherent and acyclic in degree} \(\geq m\)\uline{, there locally exists a quasi-isomorphism} \(E\longrightarrow F\)\uline{, with} \(E\) \uline{strictly} \(n\)\uline{-pseudo-coherent and} \(E^i=0\) \uline{for} \(i\geq m\)\uline{. When} \(\underline{S}\) \uline{is punctual, and} \(F\) \uline{pseudo-coherent and acyclic in degree} \(\geq m\)\uline{, there exists a quasi-isomorphism} \(E\longrightarrow F\) \uline{with} \(E\) \uline{strictly pseudo-coherent and} \(E^i=0\) \uline{for} \(i\geq m\)\uline{.}

b) \uline{Let} \(F\in\Ob\Dpx^-(\underline{C}_S)\) \uline{be acyclic in degree} \(>n\)\uline{. For} \(F\) \uline{to be} \(n\)\uline{-pseudo-coherent, it is necessary and sufficient that} \(\Ho^n(F)\) \uline{be of finite type.}

\uline{Proof.} a) Since \(F\) is acyclic in degree \(\geq m\), the arrow \(0\to F\) is an \(m\)-quasi-isomorphism. The assertion follows from (2.4) and (2.7.1).

b) There exists a distinguished triangle of \(\Dpx^-(\underline{C}_S)\):
\[
\begin{tikzcd}
& F' \arrow[dl,"+1"'] & \\
F \arrow[rr] && \Ho^n(F)[-n] \arrow[ul] & ,
\end{tikzcd}
\]
(to see this, begin by truncating \(F\), replacing \(F^n\) by \(Z^n(F)\) and \(F^i\) by \(0\) for \(i>n\)). The exact cohomology sequence implies \(\Ho^i(F')=0\) for \(i\geq n-1\), so \(F'\) is \((n-1)\)-pseudo-coherent. We conclude using (2.5) and (2.6).

% Additional notation for Expose I continuation
\providecommand{\RHom}{\mathrm R\!\Hom}
\providecommand{\Ext}{\operatorname{Ext}}
\providecommand{\coh}{\operatorname{coh}}
\providecommand{\PreCoh}{\operatorname{precoh}}
\providecommand{\Tr}{\operatorname{Tr}}
\providecommand{\Qis}{\operatorname{Qis}}
\providecommand{\Rad}{\operatorname{rad}}
\providecommand{\Spec}{\operatorname{Spec}}
\providecommand{\End}{\operatorname{End}}
\providecommand{\Id}{\operatorname{Id}}
\providecommand{\varinjlim}{\mathop{\mathrm{lim}}\limits_{\longrightarrow}}
\providecommand{\varprojlim}{\mathop{\mathrm{lim}}\limits_{\longleftarrow}}
\providecommand{\qcoh}{\operatorname{qcoh}}
\providecommand{\amp}{\operatorname{amp}}

%%%%%%%%%%  scan idx 117  |  folio 104  %%%%%%%%%%
\SGAPrintPage{104}%
\sgaStableTarget{sga6.e01.expose.occ2}%
\sgaStableTarget{sga6.e01.prop.2-11}\textbf{Proposition 2.11.} \uline{Let} \(n\in\mathbb Z\)\uline{.}

a) \uline{Let} \(F\in\Ob\Cpx^-(\underline{C}_S)\)\uline{. When, for every} \(i\geq n\) \uline{(resp.\ for every} \(i\)\uline{),} \(F^i\) \uline{is of} \((i-n)\)\uline{-finite presentation (resp.\ of} \(\infty\)\uline{-finite presentation), then} \(F\) \uline{is} \(n\)\uline{-pseudo-coherent (resp.\ pseudo-coherent).}

b) \uline{Let} \(F\in\Ob\Dpx^-(\underline{C}_S)\)\uline{. When, for every} \(i\geq n\) \uline{(resp.\ for every} \(i\)\uline{),} \(\Ho^i(F)\) \uline{is of} \((i-n)\)\uline{-finite presentation (resp.\ of} \(\infty\)\uline{-finite presentation), then} \(F\) \uline{is} \(n\)\uline{-pseudo-coherent (resp.\ pseudo-coherent).}

\uline{Proof.} It suffices to treat the assertions without the ``resp.'' alternatives.

a) We argue by induction on \(N\) such that \(F^i=0\) for \(i>n+N\). For \(N=0\), the assertion follows at once from (2.10 b)). Suppose \(N\geq1\). There is an exact sequence
\[
  0\longrightarrow F^{n+N}[-n-N]\longrightarrow F\longrightarrow F'\longrightarrow 0,
\]
By the induction hypothesis, \(F'\) is \(n\)-pseudo-coherent; on the other hand, \(F^{n+N}[-n-N]\) is \(n\)-pseudo-coherent by hypothesis (2.9 and 2.5.a)); thus, by (2.5.b)), \(F\) is \(n\)-pseudo-coherent.

b) We argue by induction on \(N\) such that \(\Ho^i(F)=0\) for \(i>n+N\). For \(N=0\), the assertion follows from (2.10 b)). Suppose \(N\geq1\). There is a distinguished triangle of \(\Dpx^-(\underline{C}_S)\) (cf.\ proof of (2.10.b))):
\[
\begin{tikzcd}
& F' \arrow[dl,"+1"'] & \\
F \arrow[rr] && \Ho^{n+N}(F)[-n-N] \arrow[ul] & .
\end{tikzcd}
\]
The exact cohomology sequence implies
\[
  \Ho^i(F')=0\quad\text{for }i\geq n+N-1,\quad\text{i.e. }i>(n-1)+(N-1),
\]
and
\[
  \Ho^i(F')\qis\Ho^{i+1}(F)\quad\text{for }i<n+N-1 .
\]
%%%%%%%%%%  scan idx 118  |  folio 105  %%%%%%%%%%
\SGAPrintPage{105}%
By the induction hypothesis (applied with \(n-1\) in place of \(n\)), \(F'\) is therefore \((n-1)\)-pseudo-coherent. On the other hand, \(\Ho^{n+N}(F)[-n-N]\) is \(n\)-pseudo-coherent (2.9 and 2.5.a)). Thus, by (2.6), \(F\) is \(n\)-pseudo-coherent.

\sgaStableTarget{sga6.e01.rem.2-11-1}\textbf{Remark 2.11.1.} When \(F\) is pseudo-coherent, the \(F^i\) (resp.\ \(\Ho^i(F)\)) are in general not pseudo-coherent, nor even of finite type. However, in the next paragraph, in the sorites on notions of coherence, we shall see important examples where the pseudo-coherence of \(F\) implies that of the \(\Ho^i(F)\).

\sgaStableTarget{sga6.e01.prop.2-12}\textbf{Proposition 2.12.} \uline{Let} \(F',F''\in\Ob\Dpx(\underline{C}_S)\)\uline{, and} \(F=F'+F''\)\uline{. For} \(F\) \uline{to be} \(n\)\uline{-pseudo-coherent (resp.\ pseudo-coherent), it is necessary and sufficient that} \(F'\) \uline{and} \(F''\) \uline{be} \(n\)\uline{-pseudo-coherent (resp.\ pseudo-coherent).}

\uline{Proof.} It suffices to prove the assertion without the ``resp.'' alternatives. Sufficiency is clear; let us prove necessity. By induction, we may suppose we have already proved that \(F'\) and \(F''\) are \((n+1)\)-pseudo-coherent. After localizing, we can even suppose that \(F'\) and \(F''\) are strictly \((n+1)\)-pseudo-coherent and that \(F=F'+F''\) in \(\Cpx^-(\underline{C}_S)\). Denote by \(F'_{>n}\) (resp.\ \(F'_{\leq n}\)) the complex obtained from \(F'\) by replacing \(F'^i\) by \(0\) for \(i\leq n\) (resp.\ \(i>n\)), and similarly \(F''_{>n}\) (resp.\ \(F''_{\leq n}\))\dots{} There is an exact sequence
\[
  0\longrightarrow F'_{>n}+F''_{>n}\longrightarrow F
  \longrightarrow F'_{\leq n}+F''_{\leq n}\longrightarrow 0 .
\]
\(F'_{>n}+F''_{>n}\) is strictly perfect and \(F\) is \(n\)-pseudo-coherent, so (2.6) \(F'_{\leq n}+F''_{\leq n}\) is \(n\)-pseudo-coherent. By (2.10 b)), this means that \(\Ho^n(F'_{\leq n}+F''_{\leq n})\) is of finite type. But we have
\[
  \Ho^n(F'_{\leq n}+F''_{\leq n})
  =\Ho^n(F'_{\leq n})+\Ho^n(F''_{\leq n}),
\]
